\فصل{جمع‌بندی و کارهای آتی}

\قسمت{جمع‌بندی پژوهش}

در این پژوهش به مسئله‌ی انطباق در زمان آزمون در شرایط برخط پرداخته شد؛ مسئله‌ای که هدف آن بهبود تعمیم‌پذیری مدل‌های از پیش آموزش‌دیده در مواجهه با تغییر توزیع داده‌ها در زمان اجرا است. در چنین شرایطی، مدل‌ها باید بدون دسترسی به برچسب داده‌ها، بدون اطلاع از کل مجموعه‌ی آزمون و بدون امکان تنظیم دقیق فراپارامترها خود را به‌صورت پویا با محیط جدید تطبیق دهند. چالش اصلی این مسئله آن است که در صورت نبود سازوکارهای کنترلی مناسب، فرآیند انطباق می‌تواند منجر به انباشت خطاو بروز پدیده‌ی فروپاشی مدل شود.  

برای درک بهتر این چالش، مروری بر روش‌های موجود انجام شد. نتایج نشان داد که بسیاری از روش‌های مبتنی بر آموزش نیازمند داده‌های منبع یا بازآموزی کامل شبکه هستند و بنابراین در محیط‌های واقعی و پویا کاربردپذیری محدودی دارند. روش‌های پسینی نیز اگرچه بدون بازآموزی عمل می‌کنند، اما عموماً بر معیارهایی مانند بیشینه‌ی خروجی یا آنتروپی تکیه دارند که در مواجهه با داده‌های خارج از توزیع دقت کافی ندارند.  

در راستای رفع این محدودیت‌ها، این پژوهش چارچوبی نوین بر پایه‌ی امتیاز جرم تقریبی معرفی نمود که بر اساس گرادیان لگاریتم درست‌نمایی نسبت به ورودی تعریف شده و قادر است داده‌های نامطمئن و خارج از توزیع را با دقت بیشتری شناسایی کند. چارچوب پیشنهادی شامل دو گام کلیدی می‌باشد:  
\begin{enumerate}
	\item پالایش اولیه‌ی داده‌ها بر اساس آنتروپی، به‌منظور حذف نمونه‌های آشکارا نامطمئن؛
	\item وزن‌دهی پویا به داده‌های باقی‌مانده بر اساس امتیاز جرم تقریبی، به‌منظور افزایش سهم نمونه‌های قابل اعتماد در فرآیند انطباق.
\end{enumerate}

این طراحی سبب می‌شود که روش پیشنهادی ضمن حفظ سادگی و عدم نیاز به بازآموزی، در هر دو حالت انطباق برخط و انطباق برخط مستمر عملکردی پایدار از خود نشان دهد. ویژگی‌های روش حاضر عبارت‌اند از:
\begin{itemize}
	\item سادگی در پیاده‌سازی با سربار کم
	\item قابلیت به‌کارگیری مستقیم بر روی مدل‌های از پیش آموزش‌دیده و انواع روش‌های انطباق زمان آزمون
	\item پایداری بیشتر در برابر تغییر توزیع داده‌ها
	\item کاهش وابستگی به تنظیم فراپارامترها و حساسیت کم‌تر
	\item قابلیت اجرا در محیط‌های بلادرنگ
\end{itemize}

بدین‌ترتیب، پژوهش حاضر به ارائه‌ی روشی ساده، پایدار و کارآمد برای انطباق در زمان آزمون می‌پردازد؛ روشی که از یک‌سو با سادگی در پیاده‌سازی و سربار محاسباتی اندک، در برابر تغییرات توزیع داده مقاوم‌تر عمل می‌کند و از سوی دیگر، توانایی تطبیق‌پذیری و انعطاف‌پذیری مدل را نیز حفظ می‌نماید.

\قسمت{کارهای آتی}

با توجه به ویژگی‌ها و محدودیت‌های روش پیشنهادی، چندین جهت برای تحقیقات آینده قابل بررسی است:

\begin{itemize}
	\item \textbf{استفاده از داده‌های منبع برای بهبود انطباق:} در صورتی که داده‌های منبع در دسترس باشند، می‌توان از روش‌های پیچیده‌تر و همچنین شناسایی لایه‌های حساس‌تر برای بهبود معیارهای مبتنی بر داده‌های خارج از توزیع استفاده کرد.  
	
	\item \textbf{گسترش به انطباق با دسته‌های کوچک:} روش حاضر برای دسته‌های نسبتا بزرگ طراحی شده است. توسعه‌ی آن برای دسته‌های کوچک و حتی تک‌نمونه‌ای نیازمند تحقیقات بیشتر و تغییرات الگوریتمی است تا در کاربردهای بلادرنگ نیز مؤثر باشد.  
	
	\item \textbf{ترکیب با روش‌های جایگزین شناسایی خارج از توزیع:} بهره‌گیری از داده‌های مصنوعی خارج از توزیع یا طراحی معیارهای جایگزین برای شناسایی داده‌های نامطمئن می‌تواند دقت و پایداری روش پیشنهادی را ارتقا دهد.  

	\item \textbf{پیاده‌سازی در حالت‌های عملیاتی:} ارزیابی روش در محیط‌های واقعی و جریان داده‌های پویا می‌تواند قابلیت تعمیم آن را بیش از پیش نشان دهد و چالش‌های عملی احتمالی را آشکار سازد.  
\end{itemize}

به طور کلی، ادامه‌ی این مسیر پژوهشی می‌تواند به توسعه‌ی روش‌هایی بینجامد که علاوه بر مقاومت در برابر تغییرات توزیع داده‌ها، در محیط‌های عملی و واقعی نیز کارایی بالایی داشته باشند و شبکه‌های عصبی عمیق را به ابزاری قابل اعتمادتر در شرایط پویا و نامطمئن بدل سازند.  
